% TOP_PROBLEM: {"id": 1282, "title": "Irrationality of Euler's Constant", "category_id": 1, "category": {"id": 1, "name": "number_theory", "display_name": "Number Theory", "description": "Properties of integers, prime numbers, Diophantine equations.", "slug": "number-theory", "order_index": 1, "created_at": "2026-07-31T15:26:25.670Z"}, "status": "open"}
% FIRST_POSTED: 2026-09-27
% !TeX program = lualatex
% ATTEMPT_STATUS: unresolved
\documentclass[11pt]{article}
\usepackage[margin=25mm]{geometry}
\usepackage{fontspec}
\setmainfont{Latin Modern Roman}
\usepackage{amsmath,amssymb,mathtools}
\usepackage{unicode-math}
\setmathfont{Latin Modern Math}
\usepackage{xurl,hyperref}
\hypersetup{colorlinks=true,urlcolor=blue}
\setcounter{secnumdepth}{0}
\setlength{\parindent}{0pt}
\setlength{\parskip}{5pt}
\setlength{\emergencystretch}{3em}
\begin{document}
\section*{\texorpdfstring{Irrationality of Euler's Constant}{Irrationality of Euler's Constant}}\label{1282}
\subsection{Definitions and mathematical statement}
Euler's constant is defined using natural logarithms by
\[
 \gamma=\lim_{n\to\infty}\left(H_n-\log n\right),
 \qquad H_n=\sum_{j=1}^n\frac1j.
\]
The selected question is whether $\gamma\notin{\mathbb{Q}}$, equivalently whether
$q\gamma-p\ne0$ for every $p\in{\mathbb{Z}}$ and integer $q\ge1$. The sequence defining the limit, rather than a finite decimal expansion, specifies the constant exactly. This is not Euler's exponential constant $e$, and transcendence of $\gamma$ would be a stronger conclusion than the requested irrationality.
\par\smallskip\textit{Formulation and concept references:} \href{https://www.ams.org/journals/bull/2013-50-04/S0273-0979-2013-01423-X/S0273-0979-2013-01423-X.pdf}{[S1]}.

\subsection{Short English statement}
Is the limiting difference between the harmonic sum and the natural logarithm an irrational number?
\subsection{Research attempt}
\paragraph{Scope and literature check.}
This attempt was prepared by GPT-6 Astra Ultra on 27 September 2026.
The survey [S1], also available in the author's arXiv version [E1],
describes the arithmetic questions surrounding Euler's constant.
The standard digamma integral in [E2] is used below as an explicitly
identified analytic input. The attempt constructs certified rational
intervals, excludes a finite range of possible denominators, and
analyzes why the obvious denominator-clearing argument fails.
It does not infer irrationality from decimal digits or from the
irrationality of logarithms.

All logarithms are natural. Set $a_n=H_n-\log n$ and
$R_n=a_n-\gamma$. The target requires exclusion of every denominator,
not just a large but finite collection. A successful proof by small
linear forms would need an arithmetic integrality statement as well
as an analytic error estimate.

\paragraph{Elementary signed remainder bounds.}
The consecutive difference is
\[
 a_k-a_{k+1}=\log(1+1/k)-\frac1{k+1}
          =\int_k^{k+1}\frac{k+1-t}{t(k+1)}\,dt>0.
\]
Summing to infinity yields $R_n>0$. On the integration interval,
\[
 \frac1{2(k+1)^2}
 <\log(1+1/k)-\frac1{k+1}
 <\frac1{2k(k+1)}.
\]
The upper sum telescopes; the lower sum exceeds the corresponding
decreasing-function integral. Therefore
\[
                         \frac1{2(n+1)}<R_n<\frac1{2n}.  \tag{1}
\]
In particular
\[
 H_n-\log n-\frac1{2n}<\gamma
             <H_n-\log n-\frac1{2(n+1)}.                \tag{2}
\]
The signs in these bounds are part of the argument. A bound on
$|R_n|$ alone would not supply the nonzero linear forms needed in an
irrationality proof.

\paragraph{A sharper integral remainder with a controlled sign.}
The digamma identities [E2], valid at every positive integer $n$, give
\[
 R_n=\frac1{2n}
       -2\int_0^\infty\frac{t}{(n^2+t^2)(e^{2\pi t}-1)}\,dt.       \tag{3}
\]
Specifically $\psi(n)=H_n-1/n-\gamma$ and the integral formula for
$\psi(n)$ imply (3). We are citing this identity, not claiming to
derive the theory of the gamma function here.

Now expand the rational factor by the finite geometric identity
\[
 \frac1{n^2+t^2}
 =\sum_{j=0}^{m-1}\frac{(-1)^jt^{2j}}{n^{2j+2}}
   +\frac{(-1)^mt^{2m}}{n^{2m}(n^2+t^2)}.               \tag{4}
\]
The remainder has a fixed sign for $t>0$. Its magnitude is strictly
less than the first omitted monomial. The moments are
\[
 2\int_0^\infty\frac{t^{2k-1}}{e^{2\pi t}-1}\,dt
       =\frac{2(2k-1)!\zeta(2k)}{(2\pi)^{2k}}
       =\frac{|B_{2k}|}{2k}.
\]
The first equality follows by the positive geometric series for the
denominator and termwise integration; the second is the classical
even-zeta evaluation in Bernoulli numbers. Thus, with
\[
 A_{n,m}=H_n-\frac1{2n}
                    +\sum_{k=1}^m\frac{B_{2k}}{2k n^{2k}},
\]
we have
\[
 \gamma=A_{n,m}-\log n+\varepsilon_{n,m},\quad
 0<(-1)^m\varepsilon_{n,m}
                <\frac{|B_{2m+2}|}{(2m+2)n^{2m+2}}.     \tag{5}
\]
For $m=4$, the correction is positive. It is important that (4) is a
finite identity with an integral remainder, rather than an assertion
that the full asymptotic series converges for fixed $n$.

\paragraph{Replacing the logarithm by a rational interval.}
For $n>1$ put $z=(n-1)/(n+1)\in(0,1)$. Integrating the geometric
series for $1/(1-z^2)$ gives
\[
 \log n=2\sum_{j=0}^{\infty}\frac{z^{2j+1}}{2j+1}.
\]
Let $L_{n,K}$ denote the sum through $j=K-1$. Its positive tail satisfies
\[
 0<\log n-L_{n,K}
      <T_{n,K}:=\frac{2z^{2K+1}}{(2K+1)(1-z^2)}.        \tag{6}
\]
Every quantity in (6), except the logarithm itself, is rational and can
be evaluated exactly. Combining (5) at $m=4$ with (6) yields
\[
 A_{n,4}-L_{n,K}-T_{n,K}<\gamma
             <A_{n,4}-L_{n,K}+\frac1{132n^{10}}.         \tag{7}
\]
Here
\[
 A_{n,4}=H_n-\frac1{2n}+\frac1{12n^2}-\frac1{120n^4}
                         +\frac1{252n^6}-\frac1{240n^8},
\]
and $B_{10}/10=1/132$ fixes the last coefficient.
The lower endpoint uses the upper bound on $\log n$; the upper
endpoint uses its lower bound. Reversing that subtraction would destroy
the certificate.

\paragraph{A finite denominator exclusion, stated at its actual scope.}
Take $n=32$, $K=400$, and denote the two rational endpoints in (7) by
$L<U$. Their width is less than $7\cdot10^{-18}$, checked by rational
cross multiplication. The local script checks every integer
$1\le q\le100000$. For each $q$, the smallest possible numerator of a
rational strictly above $L$ is
\[
                              p_q=\lfloor qL\rfloor+1.
\]
It verifies $p_q/q\ge U$ for all those $q$. Thus the certified interval
contains no rational with positive denominator at most $100000$.
Reduction to lowest terms is unnecessary: testing every denominator
also includes unreduced representations.

This is an exact finite conclusion, not a comparison with a
floating-point value of $\gamma$. The decimal width is only a readable
summary; the script uses fractions throughout. The result remains
consistent with $\gamma$ being rational with a larger denominator.
It is not asserted as a record denominator bound.

\paragraph{The simplest clearing attempt diverges.}
Let $d_n=\operatorname{lcm}(1,\ldots,n)$, so $d_nH_n$ is an integer.
One might try to multiply the defining approximation by $d_n$ and use
its small remainder. Two separate obstructions appear.
First,
\[
                 d_n(H_n-\gamma-\log n)=d_nR_n
\]
still contains $d_n\log n$, which is not an integer and has not been
arithmetically eliminated. Second, even its positive size does not
tend to zero. Since $n$ and $n-1$ are coprime and both divide $d_n$,
\[
 d_n\ge n(n-1),\qquad
 d_nR_n>\frac{n(n-1)}{2(n+1)}\longrightarrow\infty.      \tag{8}
\]
This elementary lower bound already kills that route; no prime number
theorem estimate for $d_n$ is needed.

Using more terms in (5) improves the analytic remainder, but creates
additional denominators from powers of $n$ and Bernoulli numbers.
Using (6) removes the explicit logarithm only by introducing still more
rational denominators. A valid irrationality argument must estimate
their combined growth, rather than count the number of matching digits.

\paragraph{An abstract criterion and an adversarial test.}
Suppose integers $A_j,B_j$, with $B_j>0$, could be constructed such that
\[
                      0<|B_j\gamma-A_j|\longrightarrow0.\tag{9}
\]
Then $\gamma$ would be irrational. Indeed, if $\gamma=p/q$, the
nonzero integer $B_jp-qA_j$ would eventually have absolute value less
than one. Formula (9) is stronger than saying rational approximants
$A_j/B_j$ converge quickly; it controls the error after multiplying
by the denominator and requires nonvanishing.

For a simple countertest to the weaker claim, take the rational target
$c=0$ and approximants $r_j=1/j!$. Their errors decay faster than any
fixed geometric rate, but with denominator $B_j=j!$ the nonzero
scaled error is always one. Fast convergence by itself is therefore
compatible with a rational limit. Likewise the finite exclusion above
does not construct a sequence satisfying (9).

\paragraph{Verification and final status.}
The proof of (1) is elementary and independent of the digamma input.
The signs and first omitted coefficient in (5) are checked using (4).
The interval and the denominator exclusions are reproducible exact
arithmetic. The source identity (3) and the even-zeta moments are stated
as external standard inputs with their positive-real hypotheses.

\textbf{UNRESOLVED.} The strongest arithmetic conclusion proved here
is exclusion of denominators at most $100000$ by an explicit rational
certificate, together with the obstruction (8) to naive clearing.
The first missing step toward irrationality is an infinite family of
nonzero integer linear forms as in (9), with proved denominator control.
Specific next targets are simultaneous cancellation of logarithmic
terms, denominator estimates for an accelerated integral construction,
and a nonvanishing argument surviving that cancellation. None has been
established in this attempt.

\subsection{Sources}
\begin{itemize}
\item[S1]Jeffrey C. Lagarias, "Euler's constant: Euler's work and modern developments," Bull. Amer. Math. Soc. 50 (2013), no. 4, 527-628.
\url{https://www.ams.org/journals/bull/2013-50-04/S0273-0979-2013-01423-X/S0273-0979-2013-01423-X.pdf}.
\textit{Formulation source.}
\item[E1]Jeffrey C. Lagarias, \emph{Euler's constant: Euler's work and modern developments}, author preprint of [S1].
\url{https://arxiv.org/abs/1303.1856}.
\item[E2]NIST Digital Library of Mathematical Functions, Section 5.9(ii), especially equations 5.9.15 and 5.9.16.
\url{https://dlmf.nist.gov/5.9#E15}.
\textit{Explicit analytic input for the signed integral remainder.}
\end{itemize}
\end{document}
